%! Measuring Center and Spread — Unit 2, Lesson 2.2 — Warm-Up
% ONE PAGE, blank and key. Three items at 12pt, \workrowsep 50pt, item spacing 22pt. Do not add
% a fourth item — the page is full.
% Each item is a TOOL THE NOTES PICK UP AGAIN:
%   1. Read Route 5's dot plot (spiral to Lesson 2.1) and order the values -> row 1 (median,
%      mode) and row 2 (the ordered strip the quartiles are cut from).
%   2. Add and divide -> row 1 names it the MEAN and writes it x-bar = (sum x)/n = 135/9 = 15.
%   3. -4, -2, -2, 2, 6: add (0), square and add (64), 64/4 = 16, sqrt 16 = 4 -> row 3 reveals
%      these are Route 2's DEVIATIONS: students have already computed a standard deviation.
% DATA: Route 5 = 8, 10, 12, 12, 12, 15, 18, 20, 28 (sum 135, mean 15, mode 12).
% PAGE LOCKSTEP: warmup_key differs only in -key, the header, and \blank -> \ans.
\documentclass[12pt]{article}
\usepackage{saar-article}
\usepackage{saar-boxes}
\newcommand{\TallMath}[1]{$\displaystyle #1\rule[-1.4em]{0pt}{3.2em}$}
% Word bank strip for every fill-in cluster. Defined per document; the key inherits it.
\newcommand{\wordbank}[1]{\par\vspace{2pt}\noindent\fcolorbox{goldacc}{goldbg}{\parbox{\dimexpr\linewidth-2\fboxsep-2\fboxrule\relax}{\small\textbf{Word bank:} #1}}\par\vspace{4pt}}
% Handwriting room inside every \begin{work} block. Moves the blank and the key together.
\setlength{\workrowsep}{50pt}

\begin{document}

\pageheader{Unit 2, Lesson 2.2}{Warm-Up}
% NAMESTRIP: no \namedateperiod here. The student writes their name once, on the cover.

\vspace{0.15in}

\boxguard[30]
\begin{notesbox}{Spiral Review --- Maple Grove's Route 5 Bus}
\normalsize
Nine students ride Maple Grove's Route~5 bus to school. The dot plot shows each rider's
commute, in minutes.

\begin{center}
\begin{tikzpicture}[x=0.5cm, y=0.5cm]
  \draw[thick] (4.5,0) -- (30.5,0);
  \foreach \tk in {5,...,30} { \draw (\tk,0) -- (\tk,-0.25); }
  \foreach \tk in {5,10,...,30} {
    \draw[thick] (\tk,0) -- (\tk,-0.45);
    \node[below, font=\small] at (\tk,-0.45) {\tk}; }
  \node[font=\small] at (17.5,-1.9) {commute time (minutes)};
  \foreach \p/\q in {8/1, 10/1, 12/1, 12/2, 12/3, 15/1, 18/1, 20/1, 28/1} {
    \fill[cerulean] (\p,{\q*0.75}) circle (3pt); }
\end{tikzpicture}
\end{center}

\begin{enumerate}[leftmargin=1.8em, itemsep=22pt, topsep=6pt]

  \item Read the dot plot. List the nine commute times from smallest to largest, then name
        the time that appears most often.

        \par\vspace{6pt}
        Times: \blank{9.5cm}
        \par\vspace{10pt}
        Most common: \blank{3cm}

  \item Add the nine commute times, then divide the total by $9$.

        \begin{work}
          8+10+12+12+12+15+18+20+28 &= 135 \\
          135 \div 9 &= 15 \text{ minutes}
        \end{work}

  \item Five numbers: $-4,\ -2,\ -2,\ 2,\ 6$. Add them. Then square each one, and add the
        squares.
        \wordbank{$0$ \;\textbullet\; $4$ \;\textbullet\; $16$ \;\textbullet\; $36$ \;\textbullet\; $64$}

        \vspace{2pt}
        {\centering\renewcommand{\arraystretch}{2.6}
        \begin{tabular}{@{} l C{1.3cm} C{1.3cm} C{1.3cm} C{1.3cm} C{1.3cm} C{2.2cm} @{}}
        \toprule
        \textbf{Number} & $-4$ & $-2$ & $-2$ & $2$ & $6$ & Sum: \blank{1.1cm} \\
        \midrule
        \textbf{Square} & \blank{1.1cm} & \blank{1.1cm} & \blank{1.1cm} & \blank{1.1cm} & \blank{1.1cm} & Sum: \blank{1.1cm} \\
        \bottomrule
        \end{tabular}\par}

        \vspace{8pt}
        Divide the sum of squares by $4$, then take the square root: \blank{4.2cm}

\end{enumerate}
\end{notesbox}

\end{document}
